
Por supuesto, este resultado seguiría siendo válido tal cual si, en lugar de partir de $S$-preesquemas ordinarios~$X$ e~$Y$, se partiese de preesquemas formales $\widehat{\cal{I}}$-ádicos
\marginpar{76}
$\goth{X}$ y~$\goth{Y}$ sobre~$\widehat{S}$. Permite demostrar, por ejemplo, que ciertos esquemas formales propios sobre un anillo local completo son, de hecho, algebraicos. En efecto, procediendo como en el lema~\Ref{III.4.2}, se obtiene:

\begin{corollaire}
\label{III.5.7}
Bajo las condiciones de~\Ref{III.5.6}, si~$g_0$ es un isomorfismo, lo mismo ocurre con~$\widehat{g}$.
\end{corollaire}

\noindent
(N.B.: el mismo resultado vale para las inmersiones cerradas).

Se obtiene así:

\begin{proposition}
\label{III.5.8}
Sean~$A$ un anillo local completo de ideal maximal~$\goth{m}$,
\ifthenelse{\boolean{orig}}{}{de}
cuerpo residual~$k$, y sean~$\goth{X}$ y~$\goth{Y}$ dos preesquemas formales $\goth{m}$-ádicos sobre~$A$, planos sobre~$A$ (es decir, para todo~$n$, $X_n$ e~$Y_n$ son planos sobre~$A_n=A/\goth{m}^{n+1}$). Supóngase que~$X_0=\goth{X}\otimes_Ak$ es liso sobre~$k$ y que~$\H^1(X_0,\goth{g}_{X_0/k})=0$. Entonces todo $k$-isomorfismo de~$Y_0$ sobre~$X_0$ se prolonga a un $A$-isomorfismo de~$\goth{Y}$ sobre~$\goth{X}$; esta prolongación es única si, además, $\H^0(X_0,\goth{g}_{X_0/k})=0$.
\end{proposition}

Esto da, en particular, un resultado de \emph{unicidad} para un preesquema formal liso sobre~$A$ que se reduce a un preesquema~$X_0$ dado (siempre que $\H^1(X_0,\goth{g}_{X_0/k})=0$). Además, si~$\goth{X}$ y~$\goth{Y}$ provienen de esquemas ordinarios propios sobre~$A$, sean~$X$ e~$Y$, se sabe, por el teorema de existencia de haces en geometría formal (cf. la exposición del Seminario Bourbaki, \No 182\footnote{Cf. EGA III 5.4.1 para la demostración.}), que hay una correspondencia biunívoca entre los $A$-isomorfismos $Y\isomto X$ y
los $A$-isomorfismos de los completados formales; por tanto:

\begin{corollaire}
\label{III.5.9}
El enunciado precedente~\Ref{III.5.8} sigue siendo válido sustituyendo~$\goth{X}$ y~$\goth{Y}$ por esquemas ordinarios~$X$ e~$Y$ sobre~$A$ que sean \emph{propios}.
\end{corollaire}

Por último, cuando~$\goth{X}$ es un esquema formal propio sobre~$A$ y~$\goth{Y}$ es de la forma~$\widehat{Y}$, donde~$Y$ es un esquema ordinario propio sobre~$A$, la proposición~\Ref{III.5.8} da condiciones suficientes para poder hallar un isomorfismo de~$\goth{X}$ con~$\widehat{Y}$, y por tanto para que el esquema formal~$\goth{X}$ sea, de hecho, «algebraico» (es decir, isomorfo a un~$\widehat{X}$, con~$X$ un esquema ordinario propio sobre~$A$, que quedará entonces determinado canónicamente). Esto ocurre, en particular, si~$X_0=\PP_k^r$ (o, más generalmente, si~$X_0$ es un esquema de
\ifthenelse{\boolean{orig}}{Sévéri--Brauer}{Severi--Brauer},
es decir, si se vuelve
\marginpar{77}
isomorfo al espacio proyectivo tipo sobre la clausura algebraica de~$k$): todo esquema formal propio y plano sobre~$A$, de fibra~$\PP_k^r$, es algebraizable y, más precisamente, es isomorfo al completado formal $\goth{m}$-ádico de~$\PP_A^r$. En particular (gracias al «teorema de existencia»), todo esquema ordinario propio sobre~$A$, de fibra~$\PP_k^r$, es isomorfo a~$\PP_A^r$ (siendo~$A$ un anillo local completo). Utilizando la teoría del descenso, puede demostrarse que, si~$A$ no es completo, $X$ se vuelve isomorfo a~$\PP^r$ tras una extensión finita y étale de la base $A\to A'$ (y, bajo esta forma, el resultado sigue siendo válido para una fibra que sea un esquema de
\ifthenelse{\boolean{orig}}{Sévéri--Brauer}{Severi--Brauer}).
